\documentclass[twoside]{article}
\usepackage[
  a5paper,
  margin=0.5in,
  includehead,
  headsep=4pt,
  headheight=13pt
]{geometry}
\usepackage{gregosheet}
\usepackage{gregosheet-psalm}
\usepackage{lyluatex}
\usepackage{xcolor}
\input{styles}
\input{ordo-cantus-officii/6-appendix.tex}

\defsheet{CibavitEos}
    {<-ð--0qQ0-1---1ed1-0qQ0-¨-1--3-3--1--0---1ed1eD0-ð-+-0e-eE2ed1-1ed1-1-,-0--1--3--33r-4---3rR331ed1-:-1--4--rR3-34t--4rR3E2-2-+-1--4--rR3r-3-,-34t-tT4-rR3eE2-wW1-:-12e-3eE2W1¨2e¨4rR3E2-12eE21w-ed1-.-----0--1---3--3--3--3--3--3--3--4--4--3---,---1--3--3--4--ed1-3-ws0-12ed1-?,-3-4-ed1-3-ws0-12ed1-.}
    {  Ci-bá---vit e----os   * ex á-di-pe fru-mén-----ti, al-le-----lú---ia, et de pe-tra mel-le          sa-tu-rá--vit  e------os, al-le-lú---ia, al--le--lú-----ia,   al--le---------------lú------ia. <Ps.> Ex-sul-tá-te De-o, ad-iu-tó-ri no-stro, * iu-bi-lá-te De--o Ia--cob!     E u o   u a   e.}

\defsheet[tone={8szo-a-c}]{TaplaltaOket}
    {<ôþù-2---2t--3eE2W1-2e-2---:-5-6--5--3--5--5--5--,-5---6--5----6uU6-tT45zZ5-:-3---2--3---3--5-5---2---2--,-2--2--qQ0-ñ-:-1--2e-2--2-?,}
    {     Táp-lál-ta     ő--ket * a ga-bo-na ja-vá-val, s_a kő-szik-lá---ból       jól-la-kat-ta ő-ket méz-zel, al-le-lu--ja, al-le-lu-ja.}

\defsheet[tone={8szo-a-c}]{AkiEszi}
    {<ôþù-2t-3--:-3-3---2--1--2e--3--2--:-5--6---6---5--3--5--5---5-,-6--7---6---tT45z-5--:-3--5--5-2---2-,-2--2--qQ0-ñ-:-1--2e-2--2-?,}
    {     A--ki * e-szi az én tes-te-met, és isz-sza az én vé-rem-et, én-ben-nem ma----rad, és én ő-ben-ne  al-le-lu--ja, al-le-lu-ja.}

\defsheet{MissaDeAngelisKyrie}
    {<X-3--56u-7ciI7U6¨7cpé8I7U6¨7iI7c-:-uj5T4R3¨zZ5-4--4-3y-?,-5-----tT4R3E2ey356uc-iI7U6uc-:-uj5T4R3¨zZ5-4--4-3y-?,-ö--9--pP9O8¨9pÉ7cpÉ7iv56uc-:-uj5T4R3¨zZ5-4--4-3y-?,-ö--9--pP9O8¨9pÉ7c-:-pP9¨pP9O8¨9pÉ7cpÉ7iv56uc-:-uj5T4R3¨zZ5-4--4-3y-.}
    {   Ky-ri--e                        e-----------le-i-son.  Chris-te             @         e-----------le-i-son.  Ky-ri-e                       e-----------le-i-son. Ky-ri-e             @                          e-----------le-i-son.}

\defsheet{MindenekSzemei}
    {<-0---0--1---0---0e-3-:-¢------------1---2e-4--,-¢--------------------------------1---0-1e-3-.}
    {  Min-de-nek sze-me-i   tebenned_bíz-nak U--ram, mert_te_adsz_nekik_eledelt_alkal-mas i-dő-ben.}

\begin{document}

\section*{Úrnapja}

\parttitle[ÉE 168]{Gyülekező ének}
Zálogát adtad, ó Jézus, örök szeretetednek\\
Rendelvén e nagy Szentséget és adván tieidnek.\\
Hogy teveled egyesüljünk s téged viszontszeressünk,\\
Örök hála és imádás légyen azért nevednek.\\

S hogy halálod szent emlékét mindennap megülhessük,\\
Hagytad ezt a nagy Szentséget, hogy magunkhoz vehessük.\\
Csodálatos kincset adtál, és közöttünk maradtál.\\
Véghetetlen jóságodért, add, hogy mindig szeressünk.

\parttitle[GSLH 281]{Introitus}
\printsheet{CibavitEos}

\parttitle[ÉE 575]{Introitus}
\printsheet{TaplaltaOket}

\begin{psalmsheet}[tone=8, number={80. zsoltár}]
  \psalmtext{
    Örvendezzetek a mi segítőnknek, az Istennek, * ujjongjatok Izrael urának!\\
    Kezdjetek énekbe és hozzátok elő a dobot, * a kedves lantot a citerával együtt.\\
    Én vagyok a te Urad, Istened, † aki Egyiptom földjéről kihoztalak téged, * nyisd meg szádat és én jóllakatlak téged. ANT. ---

    Dicsőség az Atyának és Fiúnak * és Szentlélek Istennek.\\
    Mikképpen kezdetben vala, most és mindenkor, * és mindörökkön örökké, amen. ANT.
}
\end{psalmsheet}

\parttitle[ÉE 474]{Kyrie}
\printsheet{MissaDeAngelisKyrie}

\parttitle[ÉE 576]{Graduále}
\printsheet{MindenekSzemei}

V.! Föltárod a te szent keze\underline{i}det, * és betöltesz minden élőket ál/d\underline{á}soddal. \textbf{TE ADSZ NEKIK... MINDENEK SZEMEI...}

\parttitle[GH 556]{Alleluja}
\printsheet{AlleluiaHypolydian}

V. Az én testem valób\underline{a}n / étel, * és az én vérem val\underline{ó}ban / ital. \textbf{ALLELUJA.}\\
V. Aki eszi az én testemet és issza az \underline{é}n / véremet, * bennem marad, \underline{é}s én / őbenne. \textbf{ALLELUJA.}

\parttitle[ÉE 557]{Szekvencia}
\begin{gregosheet}
  \addinline{<-----0--0------1--0---3---2---1--0--:-0--2--4--2----4----5--5--4--:-4---3--2---3---1-1--0-?,-------0-0-----1--0----3--2----1---0--:-0---2--4----2---4-5--5---4--:-4---3---2---3--1----1--0-?,-------4----4--5---4---7----6--5-4--:-0-2---4-2--4--5---5--4--:-4-3--2----3---1--1--0-?,--------4----4---5--4---7---6--5--4---:-0--2---4--2--4-5---5--4--:-4---3--2-3---1-----1--0-.}
            {<1.a> Di-csérd, Si-on, Meg-vál-tó-dat, ve-zé-re-det, pász-to-ro-dat, áld-ja han-gos é-ne-ked! <1.b> A-hogy, bí-rod, ak-ként mer-jed, bár-mi nagy-nak é-ne-kel-jed, mél-tón nem di-csér-he-ted! <2.a> Nagy ti-tok-ról szól az é-nek: é-let é-lő kút-fe-jé-nek, a Ke-nyér-nek hó-do-lunk, <2.b> Mely-ben mi-dőn hal-ni ké-szül, ma-gát ad-ja ö-rök-ré-szül tár-sa-i-nak Krisz-tu-sunk.}
\end{gregosheet}

\begin{hymnstanzas}
    \textcolor{red}{3.a} Legyen teljes, legyen zengő,\\
    legyen vidám s hozzá illő\\
    szívből fakadt énekünk,
\stanzabreak
    \textcolor{red}{3.b} Annak titkát fontolgatva,\\
    magát Jézus miképp adta\\
    táplálékul minekünk!
\stanzabreak
    \textcolor{red}{4.a} Az új király asztalára\\
    kenyér szállott, égi pászka,\\
    elavult a régi már.
\stanzabreak
    \textcolor{red}{4.b} Fut az újtól, ami régi,\\
    nap az éjjelt fölcseréli,\\
    fut a fénytől a homály.
\stanzabreak
    \textcolor{red}{5.a} Amit Krisztus tett ez estén,\\
    hagyta, hogy rá emlékezvén\\
    cselekedjük mi is azt.
\stanzabreak
    \textcolor{red}{5.b} Áldás fakadt törvényéből,\\
    így lesz borból és kenyérből\\
    üdvösséges áldozat.
\stanzabreak
    \textcolor{red}{6.a} Drága titka szent hitünknek:\\
    testté-vérré lényegülnek,\\
    bor s kenyér mi volt előbb.
\stanzabreak
    \textcolor{red}{6.b} Régi rend itt újnak enged,\\
    szárnya lankad észnek,\\
    szemnek, élő hitből végy erőt!
\stanzabreak
    \textcolor{red}{7.a} Más és más, de nem lényegben,\\
    csak jel szerint, más színekben\\
    mennybéli jók rejlenek.
\stanzabreak
    \textcolor{red}{7.b} Vére ital, teste étel,\\
    mégis ott van, teljességgel\\
    mindkettőben Istened.
\stanzabreak
    \textcolor{red}{8.a} Aki veszi, meg nem osztja,\\
    meg nem töri, nem szakasztja:\\
    oszthatatlan eledel.
\stanzabreak
    \textcolor{red}{8.b} Veszi egy és veszik ezren,\\
    rövidséget nem lát egy sem,\\
    fogyasztják, és nem fogy el.
\stanzabreak
    \textcolor{red}{9.a} Veszik jók, és veszik rosszak,\\
    csakhogy különféle sorsnak\\
    részesei ők vele.
\stanzabreak
    \textcolor{red}{9.b} Rossznak halál, jónak élet:\\
    ilyen az egyforma étek\\
    különböző ereje.
\end{hymnstanzas}

\begin{gregosheet}
  \addinline{<------0-0---1--0--3--2---1---0--:-0--2---4--2----4---5--5---4--,-4-4--5--4--7--6--5--4--:-4--3---2--3--1---1--0-?,-------0-0----1--0---3--2---1--0--:-0-2---4---2--4----5--5--4-,-4---4---5--4----7--6--5--4--:-4---3--2---3---1-1--0-?,-------0-0---1--0-----3---2----1----0--:-0--2----4---2--4--5--5---4--,-4--4---5-----4--7--6--5---4---:-0------2----4-2---4--5--5---4---4--3----2--3---1----1--0--?,---------0---0--1---0----3-----2--1--0---:-0---2--4---2---4---5------5--4----,-4----4----5---4-----7--6---5--4--:-0-----2---4---2-4---5---5---4---4--3--2---3---1-----1--0---?,-0qQ0-ñ]-.}
            {<10.a> Í-me, ét-ke an-gya-lok-nak, Ke-nye-re lett ván-do-rok-nak, e-le-de-le jó fi-ak-nak, eb-nek ad-ni nem va-ló. <10.b> I-zsák en-nek ké-pe, ár-nya, ál-do-zat-ra kész ol-tá-ra, hús-vét-bá-rány ál-do-zá-sa,  hul-ló man-na, é-gi jó. <11.a> I-gaz ke-nyér, leg-főbb pász-tor, Jé-zus, óvj az el-bu-kás-tól, te táp-lálj, és te pa-lás-tolj, s_add, hogy a fel-tá-ma-dás-kor le-gyen ná-lad kész he-lyünk! <11.b> Te, ki min-dent bírsz és ér-tesz, föl-di lét-ben táp-lálsz, él-tetsz, add, hogy lás-sunk, di-cső-sé-ges, s_tár-sul-ván a bol-dog nép-hez la-ko-mád-ban részt ve-gyünk! Á----men}
\end{gregosheet}

\parttitle[ÉE 155]{Felajánlási ének}
\begin{gregosheet}
  \addinline{<X-1----1---1-----0-3--3---4t-5---:-56u--6---6---5---4---6--tT4R3-,-3--4----6---5---4---3--rR3-4-:-4----5--3---2--1--4--rR3d1Q0-,-3-3----3-1---2-3---rR3-4-:-4---5---3---rR3-1---0--1--?,-1wW1-0q-.}
            {   Zeng-jed nyelv a di-cső-sé-ges * Test tit-kát s_a drá-ga vért,   me-lyet hul-lat-ván ér-té--kes vált-sá-gul az em-be-rért,     a föld U-ra, a föl-sé--ges méh gyü-möl-cse nem kí-mélt. Á----men.}
\end{gregosheet}

\begin{hymnstanzas}
  Méltó ezt a nagy szents\underline{é}get\\
  t\underline{é}rdre hullva áldan\underline{i},\\
  és a régi Szövets\underline{é}get\\
  új rítussal váltan\underline{i},\\
  pótolják a rest érz\underline{é}ket\\
  a merész h\underline{i}t szárnyai!
\stanzabreak
  Az Atyának és Fi\underline{á}nak\\
  l\underline{é}gyen áldás, dicsős\underline{é}g,\\
  üdv, hozsanna és im\underline{á}dat,\\
  ujjongások hirdess\underline{é}k,\\
  s aki Kettejükből \underline{á}rad:\\
  a Lélek \underline{i}s áldassék. \underline{Á}m\underline{e}n.
\end{hymnstanzas}

\parttitle[ÉE 477]{Sanctus}
\begin{gregosheet}
  \addinline{<X-3rtg34rR3-3y---:-eE2--qQ0í-:-34tuj5tT4R3-rR3y-;-35uU6u-zZ5-5x--7--zZ5-rR34t3ed1Q0¨¨34tuj5tT4R3-rR3-3y-,-3---5u-7c---tT4u-7c-7i-iI7U67-zZ5x-:-7---zZ5-4t-tT4R3-3y-,-3--eE2-qQ0-3--4t-5uj5tT4R3-rR3y-,-3--5u-uj5T47-7c--7i--iI7U67-zZ5c-:-5--5--3--5--77i-7--7c-,-7--7iiI756-tT4-:-4t-eE2-qq\\Q034tuj5tT4R3-rR3y-.}
            {   Sanc------tus, * sanc-tus,   sanc--------tus,   Do-----mi--nus De-us  Sa-----------------------ba--oth. Ple-ni sunt cae--li et ter----ra     glo-ri--a  tu----a.   Ho-san-na  in ex-cel-------sis.   Be-ne-dic----tus qui ve-----nit    in no-mi-ne Do--mi-ni.  Ho-san-----na    in ex--cel---------------sis.}
\end{gregosheet}

\parttitle[ÉE 478]{Agnus Dei}
\begin{gregosheet}
  \addinline{<X-3rrR3-4t--3\\r\\R3\\-3y-:-3---ed1-qQ0-1---0--1e-3\\r\\R3\\-3y-;-3--4t-5--4zZ5T4t-3\\r\\R3\\-3y-?,-3--5u--uj5T47-7c-:-7---tT4t-3---4---3--4t-3\\r\\R3\\-3y-;-3--4t-5--4zZ5T4t-3\\r\\R3\\-3y-?,-3rrR3-4t--3\\r\\R3\\-3y-:-3---ed1-qQ0-1---0--1e-3\\r\\R3\\-3y-;-3--4t-5--4zZ5T4t-3\\r\\R3\\-3y-.}
            {   Ag----nus De---------i  * qui tol-lis pec-ca-ta mun--------di,  mi-se-re-re      no---------bis.  Ag-nus De-----i  * qui tol--lis pec-ca-ta mun--------di,  mi-se-re-re      no---------bis.  Ag----nus De---------i  * qui tol-lis pec-ca-ta mun--------di,  do-na no-bis     pa---------cem.}
\end{gregosheet}

\parttitle[ÉE 187]{Áldozási ének}
Győzelemről énekeljen napekelt és napnyugat,\\
millió szív összecsengjen, magasztalja az Urat!\\
Krisztus újra földre szállott, vándorlásunk társa lett,\\
mert szerette a világot, kenyérszínbe rejtezett.\\
R) Krisztus kenyér s bor színében Úr s Kiály a föld felett,\\
forassz eggyé békességben mindennépet, nemzetet!\\

Egykor értünk testet öltött, kisgyermekként jött közénk,\\
a keresztfán vére ömlött váltságunknak béreként.\\
Most az oltár Golgotáján újra itt a drága vér,\\
áldozat az Isten-Bárány, Krisztus teste a kenyér. R.

\parttitle[GR 383]{Communio}
\begin{gregosheet}
  \addinline{<-0---1---33-3---3---3---rR3tT4R3-3y-;-3--34t-t\\T4-5u--7---7---uj5uU6Z5-5x-;-tT4-5èu-4tT4R34-tT4x-,-1--1333-q\\Q0-ed1e-rR3r-4x-;-X6zZ5zZ5T4t\\T4tg33-3---3--3\\r\\R3\\-3y-.}
            {  Qui man-du-cat car-nem me-------am * et bi--bit   san-gui-nem me-------um,  in  me  ma------net,   et e----go    in   e----o,   di------------------cit Do-mi---------nus.}
\end{gregosheet}

\parttitle[ÉE 596]{Kommunió}
\printsheet{AkiEszi}

\begin{psalmsheet}[tone=8, number={22. zsoltár}]
  \psalmtext{
    1. Az Úr az én Pásztorom, nincsen hiányom semmiben, * zöldellő réteken ád helyet nékem.\\
    2. Csöndes folyóvizek mellett nevelt föl engem, * felüdítette az én lelkemet. ANT. ---

    3. Az igazság ösvényein vezetett engem, * az ő nevéért.\\
    4. Ha a halál völgyében járok is, nem félek a rossztól, * mert te ott vagy velem. ANT. ---

    5. A te vessződ és pásztorbotod, * megvigasztal engem.\\
    6. Asztalt terítesz nekem, * hogy szorongatóim szégyent valljanak. ANT. ---

    7. Megkented olajjal fejemet, * színültig töltöd kelyhemet.\\
    8. És a te irgalmasságod kísér engem, * életemnek minden napján.\\
    9. Hogy az Úr házában lakjam, * időtlen időkig. ANT.
}
\end{psalmsheet}

\parttitle[Ho 150]{Kivonulási ének}
Uram Jézus! légy velünk, mi egyetlen örömünk!\\
Nálad nélkül keserűség itt a földön életünk.\\
Ó egyetlen örömünk, Uram Jézus, légy velünk.\\

Uram Jézus! tégedet szívünk, lelkünk úgy eped,\\
Mert kívüled nem találhat sehol biztos nyughelyet.\\
Boldog mással nem lehet, Uram Jézus, csak veled.

\end{document}